\documentclass[11pt]{article}
\usepackage[T1]{fontenc}
\usepackage{lmodern}
\usepackage[margin=1in]{geometry}
\usepackage{amsmath,amssymb,amsthm,mathtools}
\usepackage{enumitem}
\usepackage{microtype}
\usepackage{xcolor}
\PassOptionsToPackage{hyphens}{url}
\usepackage{hyperref}
\usepackage{bookmark}
\hypersetup{colorlinks=true,linkcolor=blue!45!black,citecolor=blue!45!black,
  urlcolor=blue!45!black,pdfauthor={OpenAI},
  pdftitle={A Counterexample to the Group-Ring Determinant Conjecture}}
\IfFileExists{glyphtounicode-cmex.tex}{\input{glyphtounicode-cmex.tex}}{}
\pdfgentounicode=1
\newtheorem{theorem}{Theorem}[section]
\newtheorem{proposition}[theorem]{Proposition}
\newtheorem{lemma}[theorem]{Lemma}
\theoremstyle{remark}
\newtheorem{remark}[theorem]{Remark}
\numberwithin{equation}{section}
\newcommand{\F}{\mathbb F}
\newcommand{\Z}{\mathbb Z}
\newcommand{\Q}{\mathbb Q}
\newcommand{\R}{\mathbb R}
\newcommand{\C}{\mathbb C}
\newcommand{\Mat}{\operatorname{M}}
\newcommand{\diag}{\operatorname{diag}}
\newcommand{\N}{\mathcal N}
\newcommand{\unverified}[1]{\par\smallskip\noindent\fbox{\parbox{0.95\linewidth}{\textbf{Unverified step.} #1}}\par\smallskip}
\setlength{\emergencystretch}{2em}
\title{A Counterexample to the Group-Ring Determinant Conjecture}
\author{Rewritten for readability from the OpenAI manuscript of the same title}
\date{September 23, 2026}

\begin{document}
\maketitle
\begin{abstract}
We disprove the unrestricted group-ring Determinant Conjecture. We construct a finitely generated group $G$ and a square matrix over $\Z[G]$ that is invertible over $\Q[G]$ and has Fuglede--Kadison determinant strictly between zero and one. The logarithmic integral defining the determinant is finite. The construction rests on one cited input, Theorem~\ref{thm:input}.
\end{abstract}

\section{Result, setting, and proof at a glance}
\label{sec:introduction}

A nonsingular integer matrix has determinant of absolute value at least one. The group-ring Determinant Conjecture asks for the analogue when the entries are finite integer combinations of elements of a discrete group and the determinant is that of Fuglede and Kadison~\cite{FK1952}.

\paragraph{Conventions.}
Let $G$ be a discrete group and $\Z[G]\subseteq\Q[G]\subseteq\C[G]$ its group rings of finite sums. For $a=\sum_g a_gg$ write $[1]a=a_1$ for the identity coefficient and $a^*=\sum_g\overline{a_g}\,g^{-1}$; for matrices, $(A^*)_{ij}=(A_{ji})^*$. Let $(\delta_g)_{g\in G}$ be the standard orthonormal basis of $\ell^2(G)$ and $\lambda(g)\delta_h=\delta_{gh}$ the left regular representation, that is, $(\lambda(g)\xi)(h)=\xi(g^{-1}h)$. (The source writes $L_g$; we rename it because $L_x$ later denotes a shear.) For $A\in\Mat_n(\C[G])$, the operator $T_A$ acts on column vectors $\xi\in\ell^2(G)^n$ by
\[
 (T_A\xi)_i=\sum_{j=1}^n\sum_g(A_{ij})_g\,\lambda(g)\xi_j .
\]
Since $\lambda$ is a unitary representation, $A\mapsto T_A$ is a unital $*$-homomorphism: $T_{AB}=T_AT_B$ and $T_{A^*}=T_A^*$. Let $e_i\in\ell^2(G)^n$ have $\delta_1$ in component $i$ and $0$ elsewhere. The \emph{unnormalized} trace of a bounded operator $T$ on $\ell^2(G)^n$ is
\begin{equation}\label{eq:trace}
 \tau_n(T)=\sum_{i=1}^n\langle Te_i,e_i\rangle,
 \qquad\text{so}\qquad
 \tau_n(T_A)=\sum_{i=1}^n[1]A_{ii},\quad\tau_n(I_n)=n,
\end{equation}
because $\langle\lambda(g)\delta_1,\delta_1\rangle$ is $1$ for $g=1$ and $0$ otherwise. We abbreviate $\tau_n(A)=\tau_n(T_A)$.

The probability spectral measure of the positive square is
\[
 \mu_A(\Omega)=\frac1n\sum_{i=1}^n\langle\mathbf1_\Omega(T_A^*T_A)e_i,e_i\rangle
 \quad(\Omega\subseteq[0,\infty)\text{ Borel}),
\]
and the spectral theorem gives $\tau_n(f(T_A^*T_A))=n\int f\,d\mu_A$ for bounded Borel $f$. The square-matrix instance of the conjecture asserts that for every $G$, $n$ and $A\in\Mat_n(\Z[G])$,
\begin{equation}\label{eq:conjecture}
 \int_{(0,\infty)}\log t\,d\mu_A(t)\ge0.
\end{equation}
The zero atom is omitted and the integral may be $-\infty$. The usual conjecture also permits rectangular matrices \cite[Conjecture 13.2]{Luck2002}; a square counterexample already contradicts it.

If $T_A$ is invertible, then $\|T_A^{-1}\|^{-2}I\le T_A^*T_A\le\|T_A\|^2I$, so $\mu_A$ is carried by a compact subinterval of $(0,\infty)$, on which $\log$ is bounded. With $\N(G)$ the group von Neumann algebra, the determinant convention is
\begin{equation}\label{eq:normalization}
 \log\det_{\N(G)}(T_A)=\frac12\tau_n\bigl(\log(T_A^*T_A)\bigr)
 =\frac n2\int_{(0,\infty)}\log t\,d\mu_A(t).
\end{equation}
For trivial $G$ this is $\log|\det A|$.

\paragraph{Relation to right multiplication (not used below).}
The literature lets matrices act on row vectors from the right. Let $\iota$ be the unitary $\delta_g\mapsto\delta_{g^{-1}}$ in each component. Since $\iota\lambda(g)\iota\,\delta_h=\delta_{hg^{-1}}$, for integral $A$ the operator $\iota T_A\iota$ is $\xi\mapsto\xi A^*$ (row vector times matrix). As $\iota e_i=e_i$, both operators have the same measure $\mu_A$; so \eqref{eq:conjecture} for $A$ is the right-multiplication statement for $A^*$.

\begin{theorem}\label{thm:main}
There are a finitely generated discrete group $G_D$, an integer $n\ge1$, and $A\in\Mat_n(\Z[G_D])$ invertible over $\Q[G_D]$ such that
\[
 0<\det_{\N(G_D)}(T_A)<1.
\]
The integral in Equation~\eqref{eq:conjecture} is finite and strictly negative. Thus the unrestricted group-ring Determinant Conjecture has a negative resolution.
\end{theorem}

A unital ring is \emph{directly finite} if $ab=1$ implies $ba=1$. The proof uses one external result.

\begin{theorem}[{\cite[Theorem 1.1]{CharacteristicTwo}}]
\label{thm:input}
There exist a finite field $K_2$ of characteristic two, a finitely generated group $H_2$, and finite sums $a_2,b_2\in K_2[H_2]$ with $a_2b_2=1$ and $b_2a_2\ne1$.
\end{theorem}

\unverified{Theorem~\ref{thm:input} is quoted from the companion paper and is neither proved nor checked here; the source gives no explicit $K_2,H_2,a_2,b_2$ to test. It says the group algebra $K_2[H_2]$ is not directly finite, contrary to Kaplansky's direct-finiteness conjecture, so what the present text establishes is the implication ``Theorem~\ref{thm:input} $\Rightarrow$ Theorem~\ref{thm:main}''.}

According to the source, the companion (its Section~6) specifies the field, group presentation and sums by a terminating finite prescription and uses torsion; only the existence statement is used here. The field $K_2$ is not asserted to be $\F_2$ (hence Lemma~\ref{lem:binary}), and a result only in odd characteristic would not supply the binary premise. Everything after Theorem~\ref{thm:input} is proved below.

\paragraph{Context.}
Determinant bounds belong to the theory of $L^2$-invariants~\cite{Luck2002}. The bound \eqref{eq:conjecture} is known for large classes of groups, including amenable groups \cite[Theorem 1.21, with its erratum]{Schick2001} and sofic groups \cite[Theorem 5]{ElekSzabo2005}. It is distinct from determinant \emph{approximation} by finite quotients: Kammeyer's counterexample \cite[Theorem 2 and Section 2]{Kammeyer2026} has determinant $2$ over the discrete Heisenberg group against at most $5^{1/3}$ on its congruence quotients of levels $3^i$, which does not violate the lower bound one. The commuting projections of Section~\ref{sec:projection} are related to those of \cite[Section 3]{DicksSchick2002} and \cite[Section 2]{LehnerNeuhauserWoess2008}; the analytic step is the tracial logarithmic derivative \cite[Lemma 2]{FK1952}, also in \cite[Lemma 1 and Proposition 2(d)]{deLaHarpeSkandalis1984}. Both are proved here; nothing is imported from these works.

\paragraph{Proof at a glance.}
\emph{Key idea.} Suppose a rational orthogonal projection $P$ of size $m$ has odd denominators and its reduction modulo $2$ factors through $r$ columns. Then for suitable $q=2^k$ there is an integer matrix $A$ that differs only by two block shears from an operator multiplying $r$ coordinates by $q$ and the range of $P$ by $-q^{-1}$. Shears do not change the determinant, so $\log\det T_A=(r-\tau_m(P))\log q$, which is negative once $\tau_m(P)>r$. (For trivial $G$, $\tau_m(P)=\operatorname{rank}P=\operatorname{rank}\bar P\le r$, so nothing goes wrong, as expected.) A failure of direct finiteness in characteristic two yields such a $P$ with $r=d$ and $\tau_m(P)=d+\theta$, $\theta>0$.

\begin{enumerate}[label=\arabic*.,itemsep=1pt,topsep=3pt]
\item Lemma~\ref{lem:binary}: Theorem~\ref{thm:input} gives $X,Y$ over $\F_2[H_2]$ with $XY=I_d\ne YX$.
\item Proposition~\ref{prop:separation} (uses 1): over $G_D=(\bigoplus_{h\in H_2}C_3)\rtimes H_2$ there is a projection $P$ of size $d+1$ with trace $d+\theta$ whose reduction factors through $d$ columns.
\item Lemma~\ref{lem:lift}: a factorization modulo $2$ lifts to one modulo $q=2^k$.
\item Lemma~\ref{lem:integral} (uses 3): the integral matrix $A$ and its rational inverse.
\item Lemmas~\ref{lem:logderivative} and~\ref{lem:shears}: shears preserve $\tau(\log(M^*M))$.
\item Proposition~\ref{prop:criterion} (uses 4, 5): $\log\det T_A=(r-\tau_m(P))\log q$.
\item Section~\ref{sec:assembly}: items 2 and 6 give $\det T_A=q^{-\theta}$.
\end{enumerate}
Items 3--6 hold for every group and do not depend on items 1--2.

\section{From the scalar defect to a projection}
\label{sec:projection}

All group-ring sums are finite and matrix products keep the indicated order.

\begin{lemma}\label{lem:binary}
Let $K_2$ be a finite field of characteristic two and $H_2$ a group. For $d=[K_2:\F_2]$, there is an injective unital ring homomorphism
\[
 \rho:K_2[H_2]\longrightarrow\Mat_d(\F_2[H_2]).
\]
For $a_2,b_2\in K_2[H_2]$ with $a_2b_2=1\ne b_2a_2$, their images $X,Y$ satisfy $XY=I_d$ and $C:=I_d-YX\ne0$, with
\begin{equation}\label{eq:defect}
 C^2=C,\qquad XC=0,\qquad CY=0.
\end{equation}
\end{lemma}

This turns the input into the hypothesis of Proposition~\ref{prop:separation}.

\begin{proof}
\emph{Step 1.} Choose an $\F_2$-basis of $K_2$ and let $m_\kappa\in\Mat_d(\F_2)$ be the matrix of $x\mapsto\kappa x$. Then $\kappa\mapsto m_\kappa$ is a unital ring homomorphism, injective because $m_\kappa$ applied to $1\in K_2$ returns $\kappa$.

\emph{Step 2.} Let $\rho(\sum_h\kappa_hh)$ be the matrix with $(i,j)$ entry $\sum_h(m_{\kappa_h})_{ij}\,h$. It is additive and unital. For monomials, the $(i,j)$ entry of $\rho(\kappa g)\rho(\kappa'h)$ is $\sum_l(m_\kappa)_{il}(m_{\kappa'})_{lj}\,gh=(m_{\kappa\kappa'})_{ij}\,gh$, since elements of $\F_2$ commute with group elements; this is the entry of $\rho(\kappa\kappa'gh)$. So $\rho$ is multiplicative. If $\rho(a)=0$, every $m_{\kappa_h}$ vanishes, hence every $\kappa_h=0$ by Step~1.

\emph{Step 3.} Put $X=\rho(a_2)$, $Y=\rho(b_2)$. Then $XY=\rho(1)=I_d$, and $YX=\rho(b_2a_2)\ne I_d$ by injectivity. From $(YX)^2=Y(XY)X=YX$,
\[
 C^2=I_d-2YX+(YX)^2=I_d-YX=C,\quad XC=X-(XY)X=0,\quad CY=Y-Y(XY)=0.\qedhere
\]
\end{proof}

\begin{proposition}\label{prop:separation}
Let $X,Y\in\Mat_d(\F_2[H_2])$, $d\ge1$, satisfy $XY=I_d\ne YX$. In the group $G_D=(\bigoplus_{h\in H_2}C_3)\rtimes H_2$, where $H_2$ acts by left translation on coordinate labels, there is an orthogonal projection $P\in\Mat_{d+1}(\Z[1/3][G_D])$ and a finite set $S\subseteq H_2\setminus\{1\}$ such that
\[
 \tau_{d+1}(P)=d+\theta,\qquad
 \theta=\frac13\left(\frac23\right)^{|S|}>0,\qquad\bar P=uv,
\]
where $u$ has size $(d+1)\times d$ and $v$ has size $d\times(d+1)$ over $\F_2[G_D]$. The integer $D=3^{1+|S|}$ clears all coefficients of $P$.
\end{proposition}

Here ``orthogonal projection'' means $P=P^*=P^2$, and $\bar P$ is the reduction modulo $2$ (Step~5). The proposition is combined with Proposition~\ref{prop:criterion} in Section~\ref{sec:assembly}.

\begin{proof}
\emph{Plan.} A rational idempotent $p$ built from the $C_3$ coordinates isolates one coefficient of $C=I_d-YX$. It carries the extra trace $\theta$, while modulo $2$ the defect $C$ absorbs it into a factorization through $d$ columns.

\emph{Notation.} Let $z_h$ generate the $h$-th copy of $C_3$ and $\Gamma=\bigoplus_hC_3$, an abelian normal subgroup of $G_D$. In $G_D$ we have $hz_kh^{-1}=z_{hk}$, and $H_2\le G_D$, so $\F_2[H_2]\subseteq\F_2[G_D]$. The matrix $C$ is nonzero and satisfies \eqref{eq:defect}: Step~3 of Lemma~\ref{lem:binary} used only $XY=I_d$.

\emph{Step 1 (normalizing one entry).} Choose $i,j$ and $h_0\in H_2$ such that $h_0$ has nonzero coefficient in $C_{ij}$; that coefficient is $1$, the field being $\F_2$. Hence $h_0^{-1}C_{ij}$ has identity coefficient $1$, and
\begin{equation}\label{eq:support}
 h_0^{-1}C_{ij}=1+\sum_{s\in S}s
\end{equation}
for a finite $S\subseteq H_2\setminus\{1\}$, possibly empty.

\emph{Step 2 (the idempotent $p$).} Put
\[
 p_h=\frac{1+z_h+z_h^2}{3},\qquad p=p_1\prod_{s\in S}(1-p_s).
\]
Since $z_h^3=1$, $(1+z_h+z_h^2)^2=3(1+z_h+z_h^2)$, and $z_h^*=z_h^2$; so $p_h=p_h^*=p_h^2$. All $p_h$ lie in the commutative ring $\Q[\Gamma]$. Hence each $1-p_s$ is a self-adjoint idempotent, and so is the commuting product $p$. As $3p_1$ and $3(1-p_s)$ are integral, $3^{1+|S|}p$ is integral. (Compare \cite[Lemmas 3.1 and 3.4]{DicksSchick2002}.)

\emph{Step 3 ($psp=0$ for $s\in S$).} From $sz_hs^{-1}=z_{sh}$ we get $sp_hs^{-1}=p_{sh}$, so
\[
 sps^{-1}=p_s\prod_{t\in S}(1-p_{st}).
\]
This contains the factor $p_s$, while $p$ contains $1-p_s$. All factors commute and $(1-p_s)p_s=0$, so $p\,(sps^{-1})=0$ and
\begin{equation}\label{eq:separation}
 psp=p(sps^{-1})s=0.
\end{equation}
Only these two factors are used; the label sets $\{1\}\cup S$ and $s(\{1\}\cup S)$ may otherwise overlap, and inverse pairs or finite-order elements in $S$ change nothing.

\emph{Step 4 ($[1]p=\theta$).} The factors $p_1$ and $1-p_s$ ($s\in S$) are supported in the coordinates $1$ and $s$, which are pairwise distinct because $1\notin S$. A product of elements of distinct coordinates of $\Gamma$ is the identity only if every factor is. So the identity coefficient of $p$ is the product of those of its factors, namely $[1]p_1=\frac13$ and $[1](1-p_s)=\frac23$. By \eqref{eq:trace},
\begin{equation}\label{eq:theta}
 \tau_1(p)=[1]p=\frac13\prod_{s\in S}\frac23=\theta>0.
\end{equation}

\emph{Step 5 (reduction modulo two).} Since $3$ is invertible in $\F_2$, there is a ring homomorphism $\Z[1/3]\to\F_2$; applied to coefficients it gives a ring homomorphism $a\mapsto\bar a$ into $\F_2[G_D]$, and likewise on matrices. Thus $\bar p^2=\bar p$ and $\bar ps\bar p=0$ by \eqref{eq:separation}. Multiplying \eqref{eq:support} by $\bar p$ on both sides,
\[
 \bar p\,h_0^{-1}C_{ij}\,\bar p=\bar p^2+\sum_{s\in S}\bar ps\bar p=\bar p.
\]

\emph{Step 6 (factorization).} Let $f$ be the row of length $d$ with entry $\bar ph_0^{-1}$ at position $i$ and zeros elsewhere, and $w$ the column with entry $\bar p$ at position $j$ and zeros elsewhere. Then $fCw=\bar ph_0^{-1}C_{ij}\bar p=\bar p$. Define
\[
 P=\diag(I_d,p),\qquad
 u=\begin{pmatrix}X\\fC\end{pmatrix},\qquad
 v=\begin{pmatrix}Y&Cw\end{pmatrix}.
\]
Using $XY=I_d$, $XC=0$, $CY=0$, $C^2=C$ in turn,
\[
 uv=\begin{pmatrix}XY&XCw\\fCY&fC^2w\end{pmatrix}
 =\begin{pmatrix}I_d&0\\0&fCw\end{pmatrix}=\diag(I_d,\bar p)=\bar P.
\]
By Step~2, $P=P^*=P^2$ and $D=3^{1+|S|}$ clears its coefficients; by \eqref{eq:trace} and \eqref{eq:theta}, $\tau_{d+1}(P)=d+\theta$.
\end{proof}

\section{Lifting the factorization and preserving integrality}
\label{sec:lifting}

From here to the end of Section~\ref{sec:analysis}, $G$ is any group. If $D>0$ is odd and $k\ge1$, then $D$ is a unit in $\Z/2^k\Z$, so reduction of integers extends to a ring homomorphism $\Z[1/D]\to\Z/2^k\Z$. Applied to coefficients it is a ring homomorphism on matrices over $\Z[1/D][G]$; for $k=1$ it is written $\bar{\ }$.

\paragraph{Why the matrix below.}
For $q=2^k>1$ and a projection $P$ of size $m$ put
\[
 J=-q^{-1}P-(I_m-P),\qquad B=\diag(qI_r,J).
\]
The absolute value of $B$ expands $r$ coordinates by $q$ and contracts the range of $P$ by $q^{-1}$, which suggests $\log\det=(r-\tau_m(P))\log q$; but $B$ is not integral. A block matrix with upper-left block $qI_r$ and off-diagonal blocks $V$ (upper), $U$ (lower) has Schur complement equal to its lower-right block minus $UV/q$. For this to be $J$, the lower-right block must be $UV/q+J=q^{-1}\bigl(UV-P-q(I_m-P)\bigr)$, which is integral exactly when $UV\equiv P+q(I_m-P)$ modulo $q$. The minus signs in $J$ produce the difference $UV-P$; Lemma~\ref{lem:lift} supplies $U,V$.

\begin{lemma}\label{lem:lift}
Let $G$ be a group, $m,r\ge1$, $D$ an odd positive integer, and $P\in\Mat_m(\Q[G])$ satisfy $P^2=P$ and $DP\in\Mat_m(\Z[G])$. Suppose $\bar P=uv$ over $\F_2[G]$, where $u$ and $v$ have respective sizes $m\times r$ and $r\times m$. For $q=2^k$, $k\ge1$, the reduction $E$ of $P$ modulo $q$ has a factorization $E=u_qv_q$ of the same sizes over $(\Z/q\Z)[G]$.
\end{lemma}

(The source calls this reduction $e=P_q$; we write $E$ because $e_i$ are basis vectors.) The lemma feeds Lemma~\ref{lem:integral}.

\begin{proof}
\emph{Plan.} Lift $u,v$ arbitrarily; the lifted product is correct modulo $2$ inside the corner cut out by $E$, and the error is removed by a finite geometric series.

\emph{Step 1.} Put $\mathcal R=(\Z/q\Z)[G]$. Reduction is a ring homomorphism, so $E^2=E$; and the image of $E$ under $\mathcal R\to\F_2[G]$ is $\bar P$.

\emph{Step 2.} Choose lifts $u_0,v_0$ of $u,v$ over $\mathcal R$, and set
\[
 Z=Eu_0v_0E,\qquad W=E-Z.
\]
Both lie in the corner ring $E\Mat_m(\mathcal R)E$, whose identity is $E$: $EW=W=WE$.

\emph{Step 3 ($W^k=0$).} Modulo $2$, $Z$ becomes $\bar Puv\bar P=\bar P^3=\bar P$, the image of $E$. So every coefficient of every entry of $W$ is even. An entry of $W^k$ is a sum of products of $k$ entries of $W$, so its coefficients are divisible by $2^k=q$ and vanish in $\Z/q\Z$.

\emph{Step 4 (corner inverse).} Let $Q=E+W+\cdots+W^{k-1}$ (so $Q=E$ if $k=1$). Since $QE=Q$ and $Z=E-W$, telescoping gives $QZ=Q-QW=E-W^k=E$.

\emph{Step 5.} Set $u_q=QEu_0$ and $v_q=v_0E$. Then $u_qv_q=Q(Eu_0v_0E)=QZ=E$. Only finite sums and products occurred, so supports stay finite.
\end{proof}

\paragraph{Choice of $q$.}
Since $2$ is a unit modulo the odd integer $D$, some $k\ge1$ has $2^k\equiv1\pmod D$ (any $k\ge1$ if $D=1$). Fix $q=2^k>1$ with $q\equiv1\pmod D$.

\begin{lemma}\label{lem:integral}
Suppose $P=P^*=P^2\in\Mat_m(\Q[G])$, an odd $D>0$ clears its denominators, and $\bar P=uv$ factors through $r\ge1$ columns. Let $q=2^k>1$ satisfy $q\equiv1\pmod D$. Choose arbitrary lifts $U,V$ over $\Z[G]$ of the factors $u_q,v_q$ supplied by Lemma~\ref{lem:lift}, so that $UV$ reduces to $E$ modulo $q$. Then
\begin{equation}\label{eq:integral-matrix}
 A=\begin{pmatrix}
 qI_r&V\\
 U&q^{-1}\bigl(UV-P-q(I_m-P)\bigr)
 \end{pmatrix}
 \in\Mat_{r+m}(\Z[G])
\end{equation}
has an inverse in $\Mat_{r+m}(\Q[G])$.
\end{lemma}

This is the matrix of Theorem~\ref{thm:main}; Equation~\eqref{eq:path} below is reused in Proposition~\ref{prop:criterion}.

\begin{proof}
\emph{Step 1 (integrality).} Write $P+q(I_m-P)=qI_m+\frac{1-q}{D}\,DP$; it is integral because $D\mid q-1$ and $DP$ is integral. Its reduction modulo $q$ is $0+(1-0)E=E$, which is also the reduction of $UV$. So every coefficient of $UV-P-q(I_m-P)$ is an integer divisible by $q$, and $A$ is integral.

\emph{Step 2 (shear factorization).} With $J,B$ as above, $P(I_m-P)=0$ gives $J\bigl(-qP-(I_m-P)\bigr)=P+(I_m-P)=I_m$, so $J^{-1}=-qP-(I_m-P)$ and $B^{-1}=\diag(q^{-1}I_r,J^{-1})$. For $x\in[0,1]$ put
\[
 L_x=\begin{pmatrix}I_r&0\\xU/q&I_m\end{pmatrix},\qquad
 R_x=\begin{pmatrix}I_r&xV/q\\0&I_m\end{pmatrix}.
\]
Then $L_xB=\begin{pmatrix}qI_r&0\\xU&J\end{pmatrix}$, and multiplying by $R_x$ on the right,
\begin{equation}\label{eq:path}
 M(x):=L_xBR_x=
 \begin{pmatrix}qI_r&xV\\xU&x^2UV/q+J\end{pmatrix}.
\end{equation}
Since $UV/q+J=q^{-1}(UV-P-q(I_m-P))$, we get $A=M(1)$.

\emph{Step 3 (inverse).} The inverses of $L_x,R_x$ are obtained by negating the off-diagonal block. Hence $M(x)^{-1}=R_x^{-1}B^{-1}L_x^{-1}$, whose entries are polynomials in $x$ with coefficients in $\Q[G]$; at $x=1$ this inverts $A$ over $\Q[G]$. In particular $x\mapsto T_{M(x)}$ is a polynomial, hence continuously differentiable, path of invertible operators. No identities $PU=U$, $VP=V$ or $V=U^*$ are required of the lifts.
\end{proof}

\section{The trace calculation}
\label{sec:analysis}

Let $\mathcal A_n$ be the operator-norm closure of $\{T_a:a\in\Mat_n(\C[G])\}$, a unital $C^*$-algebra. Three facts about $\tau_n$ from \eqref{eq:trace} are used.
\begin{enumerate}[label=(T\arabic*),itemsep=1pt,topsep=3pt]
\item $|\tau_n(T)|\le n\|T\|$ and $\tau_n(T^*)=\overline{\tau_n(T)}$, directly from $\tau_n(T)=\sum_i\langle Te_i,e_i\rangle$.
\item $\tau_n(ab)=\tau_n(ba)$ for $a,b\in\Mat_n(\C[G])$: both sides equal $\sum_{i,j}\sum_g(a_{ij})_g(b_{ji})_{g^{-1}}$.
\item By (T1), $\tau_n$ is norm-continuous, so (T2) extends from the dense subalgebra to all of $\mathcal A_n$.
\end{enumerate}

\begin{lemma}\label{lem:logderivative}
For a continuously differentiable path $M:[0,1]\to\mathcal A_n$ of invertible elements,
\[
 \frac{d}{dx}\tau_n\bigl(\log(M(x)^*M(x))\bigr)
 =2\operatorname{Re}\tau_n\bigl(M'(x)M(x)^{-1}\bigr).
\]
\end{lemma}

This is \cite[Lemma 2]{FK1952} with the unnormalized trace; it is used only in Lemma~\ref{lem:shears}.

\begin{proof}
\emph{Plan.} Expand $\log N$, $N=M^*M$, in a norm-convergent power series, differentiate termwise, and use cyclicity to reach $\tau_n(N^{-1}N')$.

\emph{Step 1 (bounds).} Inversion is continuous, so $\|M(x)\|$ and $\|M(x)^{-1}\|$ are bounded on $[0,1]$. Since $\langle N\xi,\xi\rangle=\|M\xi\|^2\ge\|M^{-1}\|^{-2}\|\xi\|^2$, there are constants $0<\delta<c$ with $\delta I\le N(x)\le cI$ for all $x$.

\emph{Step 2 (series).} Put $\Phi=I-N/c$ (the source's $W$). Then $0\le\Phi\le\eta I$ with $\eta=1-\delta/c<1$, so $\|\Phi\|\le\eta$, and functional calculus applied to $N=c(I-\Phi)$ gives, with norm convergence uniform in $x$,
\[
 \log N=(\log c)I-\sum_{j\ge1}\frac{\Phi^j}{j}.
\]

\emph{Step 3 (termwise differentiation).} Here $N'=M'^*M+M^*M'$ and $\Phi'=-N'/c$ are continuous. The product rule gives $(\Phi^j)'=\sum_{i=0}^{j-1}\Phi^i\Phi'\Phi^{j-1-i}$, hence
\[
 \left\|\frac{d}{dx}\frac{\Phi^j}{j}\right\|\le\eta^{j-1}\sup_{y\in[0,1]}\|\Phi'(y)\|.
\]
The series of derivatives therefore converges uniformly, so $\log N$ is continuously differentiable with termwise derivative. By (T1), $\tau_n$ commutes with the sum and the derivative.

\emph{Step 4 (cyclicity).} By (T3) each of the $j$ summands of $(\Phi^j)'$ has trace $\tau_n(\Phi^{j-1}\Phi')$; commutation of $N$ and $N'$ is not assumed. Therefore
\[
 \frac{d}{dx}\tau_n(\log N)=-\sum_{j\ge1}\tau_n(\Phi^{j-1}\Phi')
 =\frac1c\,\tau_n\Bigl(\sum_{j\ge1}\Phi^{j-1}N'\Bigr)=\tau_n(N^{-1}N'),
\]
because $\sum_{j\ge1}\Phi^{j-1}=(I-\Phi)^{-1}=cN^{-1}$.

\emph{Step 5 (two summands).} Use $N^{-1}=M^{-1}(M^*)^{-1}$. First, $\tau_n(N^{-1}M^*M')=\tau_n(M^{-1}M')=\tau_n(M'M^{-1})$. Second, by cyclicity and (T1),
\[
 \tau_n(N^{-1}M'^*M)=\tau_n(MN^{-1}M'^*)=\tau_n\bigl((M^*)^{-1}M'^*\bigr)
 =\tau_n\bigl((M'M^{-1})^*\bigr)=\overline{\tau_n(M'M^{-1})}.
\]
The sum is $2\operatorname{Re}\tau_n(M'M^{-1})$.
\end{proof}

\begin{lemma}\label{lem:shears}
Let $r,m\ge1$ and let $B\in\mathcal A_{r+m}$ be invertible. Let $\alpha,\beta$ have respective sizes $m\times r$ and $r\times m$ over the operator-norm closure of $\C[G]$. For $x\in[0,1]$, let $L_x,R_x$ be lower and upper block shears with diagonal blocks $I_r,I_m$ and off-diagonal blocks $x\alpha,x\beta$, respectively. Then $\tau_{r+m}(\log((L_xBR_x)^*(L_xBR_x)))$ is constant.
\end{lemma}

(The source calls the blocks $S,T$; we rename them since $S$ is the set of Proposition~\ref{prop:separation}.) Used in Proposition~\ref{prop:criterion} with $\alpha=U/q$, $\beta=V/q$.

\begin{proof}
Write $M=L_xBR_x$, a polynomial path with $M^{-1}=R_x^{-1}B^{-1}L_x^{-1}$, the shear inverses negating the off-diagonal blocks. From $M'=L_x'BR_x+L_xBR_x'$,
\[
 M'M^{-1}=L_x'L_x^{-1}+(L_xB)(R_x'R_x^{-1})(L_xB)^{-1}.
\]
Direct multiplication gives
\[
 L_x'L_x^{-1}=\begin{pmatrix}0&0\\\alpha&0\end{pmatrix},\qquad
 R_x'R_x^{-1}=\begin{pmatrix}0&\beta\\0&0\end{pmatrix}.
\]
Both have zero diagonal entries, hence trace zero by \eqref{eq:trace}. By (T3) the conjugate $(L_xB)(R_x'R_x^{-1})(L_xB)^{-1}$ also has trace zero. So $\tau_{r+m}(M'M^{-1})=0$, and Lemma~\ref{lem:logderivative} gives zero derivative.
\end{proof}

\begin{proposition}\label{prop:criterion}
Let $G$ be a discrete group, $m,r\ge1$, and let $P=P^*=P^2\in\Mat_m(\Q[G])$. Suppose an odd $D>0$ clears its denominators and $\bar P=uv$ with $u\in\Mat_{m,r}(\F_2[G])$ and $v\in\Mat_{r,m}(\F_2[G])$. For every $q=2^k>1$ congruent to one modulo $D$, the matrix $A$ of Lemma~\ref{lem:integral} satisfies
\begin{align}
 \tau_{r+m}(\log(T_A^*T_A))
   &=2(r-\tau_m(P))\log q,\label{eq:criterion-trace}\\
 \log\det_{\N(G)}(T_A)&=(r-\tau_m(P))\log q.
 \label{eq:criterion-det}
\end{align}
It is integral and rationally invertible, and its positive-square spectral measure has support in a compact subinterval of $(0,\infty)$.
\end{proposition}

\begin{proof}
\emph{Step 1.} Integrality and the inverse over $\Q[G]$ are Lemma~\ref{lem:integral}. Hence $T_A$ has the bounded inverse $T_{A^{-1}}$, and $\|T_A^{-1}\|^{-2}I\le T_A^*T_A\le\|T_A\|^2I$ gives the support statement.

\emph{Step 2 (shear invariance).} Apply Lemma~\ref{lem:shears} to $T_B$ (invertible by Step~2 of Lemma~\ref{lem:integral}) with $\alpha=T_{U/q}$, $\beta=T_{V/q}$. By \eqref{eq:path} the path is $T_{M(x)}$, with $M(1)=A$ and $M(0)=B$. So $\tau_{r+m}(\log(T_A^*T_A))=\tau_{r+m}(\log(T_B^*T_B))$.

\emph{Step 3 (value at $x=0$).} Since $P=P^*=P^2$, $J^*=J$ and $J^2=q^{-2}P+(I_m-P)$, so
\[
 B^*B=\diag\bigl(q^2I_r,\ q^{-2}P+I_m-P\bigr).
\]
For a projection $P$ and real $s$, $\exp(sP)=I_m+(e^s-1)P$; with $s=-2\log q$ this is $q^{-2}P+I_m-P$. Hence
\[
 \log(B^*B)=2\log q\,\diag(I_r,-P),
\]
with trace $2(r-\tau_m(P))\log q$. This proves \eqref{eq:criterion-trace}; half of it is \eqref{eq:criterion-det} by \eqref{eq:normalization}.

\emph{Step 4.} Dividing \eqref{eq:criterion-trace} by $r+m$, as in \eqref{eq:normalization},
\begin{equation}\label{eq:criterion-integral}
 \int_{(0,\infty)}\log t\,d\mu_A(t)=\frac{2(r-\tau_m(P))\log q}{r+m}.\qedhere
\end{equation}
\end{proof}

\section{Assembly and exact determinant}
\label{sec:assembly}

\begin{proof}[Proof of Theorem~\ref{thm:main}]
Take $K_2,H_2,a_2,b_2$ from Theorem~\ref{thm:input} and put $d=[K_2:\F_2]\ge1$. Lemma~\ref{lem:binary} gives $X,Y\in\Mat_d(\F_2[H_2])$ with $XY=I_d\ne YX$. Proposition~\ref{prop:separation} then gives $G_D=(\bigoplus_{H_2}C_3)\rtimes H_2$ and an orthogonal projection $P$ of size $d+1$ with
\[
 \tau_{d+1}(P)=d+\theta,\qquad\theta=\frac13\left(\frac23\right)^{|S|}>0,\qquad D=3^{1+|S|},
\]
where the odd integer $D$ clears the coefficients of $P$ and $\bar P$ factors through $d$ columns. Choose $k\ge1$ with $q=2^k\equiv1\pmod D$. Proposition~\ref{prop:criterion} with $r=d$, $m=d+1$ produces $A\in\Mat_{2d+1}(\Z[G_D])$, invertible over $\Q[G_D]$, with
\[
 \log\det_{\N(G_D)}(T_A)=\bigl(d-(d+\theta)\bigr)\log q=-\theta\log q.
\]
Since $q>1$ and $\theta>0$,
\begin{equation}\label{eq:final-det}
 0<\det_{\N(G_D)}(T_A)=q^{-\theta}<1.
\end{equation}
With $n=2d+1$, Equation~\eqref{eq:criterion-integral} reads
\begin{equation}\label{eq:final-integral}
 \int_{(0,\infty)}\log t\,d\mu_A(t)=-\frac{2\theta\log q}{2d+1}<0,
\end{equation}
finite by the support statement of Proposition~\ref{prop:criterion}. Finally, $G_D$ is generated by a finite generating set of $H_2$ together with $z_1$, since $z_h=hz_1h^{-1}$.
\end{proof}

\paragraph{Remarks.}
The group $G_D$ differs from the companion's $H_2$ and contains elements of order three. The spectrum of $T_A^*T_A$ avoids zero, so this is a failure of the lower bound, not of determinant class. The source adds an external consequence, not used above: since $T_{A^*A}=T_A^*T_A$ is positive, \eqref{eq:normalization} gives $\log\det T_{A^*A}=\tau_n(\log(T_A^*T_A))=-2\theta\log q<0$, so applying the sofic determinant bound \cite[Theorem 5]{ElekSzabo2005} to the positive integral matrix $A^*A$ shows that $G_D$ is not sofic.

\section{Verification checklist}
\label{sec:checklist}

\begin{enumerate}[label=\arabic*.,itemsep=2pt,topsep=3pt]
\item \textbf{The external input.} Theorem~\ref{thm:input} is the only ingredient not proved here, and everything is conditional on it (see the boxed note).
\item \textbf{The trace gap.} $psp=0$ and $[1]p=\theta$: Proposition~\ref{prop:separation}, Steps~3--4. Check that $sps^{-1}$ contains $p_s$ and that the labels $\{1\}\cup S$ are distinct.
\item \textbf{Factorization through $d$ columns.} $\bar ph_0^{-1}C_{ij}\bar p=\bar p$ and $uv=\bar P$: Proposition~\ref{prop:separation}, Steps~5--6, using all three identities in \eqref{eq:defect}.
\item \textbf{Lifting.} $W^k=0$ and $QZ=E$: Lemma~\ref{lem:lift}, Steps~3--4.
\item \textbf{Integrality of $A$.} Lemma~\ref{lem:integral}, Step~1; it needs both $q\equiv1\pmod D$ and $UV\equiv E\pmod q$.
\item \textbf{$A=L_1BR_1$.} The block product \eqref{eq:path}.
\item \textbf{Shear invariance.} Termwise differentiation and cyclicity (Lemma~\ref{lem:logderivative}, Steps~3--5); vanishing diagonal of $L_x'L_x^{-1}$, $R_x'R_x^{-1}$ (Lemma~\ref{lem:shears}).
\item \textbf{Final value and sign.} $\log(B^*B)=2\log q\,\diag(I_r,-P)$ (Proposition~\ref{prop:criterion}, Step~3), then $r-\tau_m(P)=-\theta$ (Section~\ref{sec:assembly}).
\end{enumerate}

\begingroup
\raggedright
\bibliographystyle{plain}
\bibliography{references}
\endgroup
\end{document}
